\begin{titlepage}
\centering

\topskip0pt
\vspace*{\fill}
  {\Huge Creating Reproducible Environments with Nix for Scientific Computing}
\newline
\newline
Morphy Kuffour 

\href{mailto:morphy.kuffour@uconn.edu}{morphy.kuffour@uconn.edu}

Department of Computer Science and Engingeering, University of Connecticut

Thesis Advisor: Prof. Clay Tabor

Honors Advisor: Prof. Caiwen Ding
\vspace*{\fill}

\afterpage{\blankpage}
\end{titlepage}

\topskip0pt
\vspace*{\fill}
\begin{center}
  {\large \textbf{Abstract}}
\end{center}
The Nix package manager provides a way to manage dependencies between multiple software packages and configurations, enabling users to access various versions of a package and its dependencies. This allows users to maintain reproducible environments for the Community Earth System Model Version 2.0 (CESM 2.0) code across different platforms, specifically for its computational components, enabling developers to have a consistent development environment for scientific computing. Managed environments also provide a way for users to easily access and install new packages specific to the CESM project, as well as quickly update its existing packages. Using the nix package manager for reproducible development environments helps ensure accurate and consistent CESM project results, and makes it easier for developers to collaborate on the same project.
\vspace*{\fill}

\pagebreak
